% Yuvion single-column technical report, based on official NeurIPS 2026.
% Main document for Overleaf. Compile with pdfLaTeX + BibTeX.
\documentclass[10pt,letterpaper]{article}
\PassOptionsToPackage{numbers,sort&compress}{natbib}
\usepackage[preprint]{neurips_2026}
\usepackage[utf8]{inputenc}
\usepackage[T1]{fontenc}
\usepackage{microtype}
\usepackage{amsmath,amssymb}
\usepackage{booktabs}
\usepackage{tabularx}
\usepackage{algorithm}
\usepackage{algorithmic}
\usepackage{yuvion-neurips}
\usepackage[breaklinks,colorlinks,allcolors=yuvionblue]{hyperref}
\usepackage[capitalize,noabbrev]{cleveref}
\hypersetup{pdftitle={Yuvion Team Paper Template and Author Guidelines},
  pdfauthor={Yuvion Team, Alibaba Group}}

\title{Yuvion Team Paper Template\\and Author Guidelines}
\author{First Author\\
  Yuvion Team, Alibaba Group\\
  \texttt{first.author@example.com}
  \And
  Second Author\\
  Yuvion Team, Alibaba Group\\
  \texttt{second.author@example.com}}

\begin{document}
\maketitle
\pagestyle{yuvionheader}
\thispagestyle{yuvionheader}

\begin{abstract}
This template provides a single-column format for technical reports by the
Yuvion Team at Alibaba Group. It combines the typography, title treatment,
section hierarchy, and page geometry of the official NeurIPS 2026 style with
the Alibaba Security and Yuvion branded header from the team's CVPR template.
Replace this paragraph with a concise statement of the research problem,
method, principal findings, and limitations. The examples demonstrate
equations, figures, tables, algorithms, citations, and an appendix without
introducing external figure dependencies.
\end{abstract}

\section{Using this template}
\label{sec:intro}
Open \texttt{yuvion\_team\_neurips.tex} as the main document. Replace the
title, author names, email addresses, and example content. Use
\textbf{Yuvion Team, Alibaba Group} as the team affiliation. All manuscript
text should be written in English.

This is a branded technical-report template. The original CVPR template
remains available for two-column reports. For an actual conference
submission, use the conference's official author kit and its current
anonymization, page-limit, and checklist requirements.

\subsection{Single-column layout}
The official style defines a $5.5$-inch-wide by $9$-inch-high text area on
US Letter paper, with 10-point body text and 11-point leading. Keep figures,
tables, and equations within \verb|\linewidth|. Do not load the CVPR style,
enable \texttt{twocolumn}, or override the body margins with another layout
package. The branding layer adjusts only the header clearance and notice.

\subsection{Authors and references}
Separate author blocks with \verb|\And|; use \verb|\AND| to force another
author row. The bibliography uses numeric citations. For example, cite the
original Transformer paper as \citep{vaswani2017attention}, or refer to its
authors with \citet{vaswani2017attention}. Replace the example bibliography
with references actually used in your manuscript.

\section{Method and mathematical notation}
\label{sec:method}
Define each symbol when it first appears. For a model $f_\theta$ with
parameters $\theta$ and a dataset $\mathcal{D}$, an example objective is
\begin{equation}
  \theta^* = \operatorname*{arg\,min}_{\theta}
  \frac{1}{|\mathcal{D}|}\sum_{(x,y)\in\mathcal{D}}
  \ell\bigl(f_\theta(x),y\bigr).
  \label{eq:objective}
\end{equation}
Use numbered equations for expressions referenced in the text.
\Cref{eq:objective} illustrates automatic cross-references; the same commands
work for \cref{sec:method,fig:pipeline,tab:example,alg:example}.

\begin{figure}[t]
  \centering
  \setlength{\fboxsep}{10pt}
  \fbox{\parbox[c][0.42in][c]{0.23\linewidth}{\centering Input}}
  \hfill$\longrightarrow$\hfill
  \fbox{\parbox[c][0.42in][c]{0.23\linewidth}{\centering Method}}
  \hfill$\longrightarrow$\hfill
  \fbox{\parbox[c][0.42in][c]{0.23\linewidth}{\centering Output}}
  \caption{An example full-width pipeline. Replace these placeholder boxes
  with your method figure. Captions appear below figures.}
  \label{fig:pipeline}
\end{figure}

\section{Experiments and presentation}
\label{sec:experiments}
Describe the datasets, baselines, metrics, implementation details, and
evaluation protocol before reporting results. Distinguish measured results
from hypotheses and illustrations. The table below intentionally contains
no invented experimental values.

\begin{table}[ht]
  \caption{Example results table. Dashes are placeholders, not measurements.}
  \label{tab:example}
  \centering
  \begin{tabular*}{\linewidth}{@{\extracolsep{\fill}}lccc@{}}
    \toprule
    Method & Quality $\uparrow$ & Cost $\downarrow$ & Latency $\downarrow$\\
    \midrule
    Baseline & -- & -- & --\\
    Proposed method & -- & -- & --\\
    Ablation & -- & -- & --\\
    \bottomrule
  \end{tabular*}
\end{table}

\subsection{Figures, tables, and algorithms}
Use \verb|figure| and \verb|table| for single-column floats; starred floats
are not necessary. Put table captions above the table and figure captions
below the figure. Use legible labels and avoid encoding distinctions only
by color. Insert a graphic with:
\begin{verbatim}
\includegraphics[width=\linewidth]{figure.pdf}
\end{verbatim}

\begin{algorithm}[ht]
  \caption{Example evaluation procedure}
  \label{alg:example}
  \begin{algorithmic}[1]
    \REQUIRE Evaluation set $\mathcal{D}$, fixed model $f_\theta$
    \STATE Initialize results $\mathcal{R}\gets\varnothing$
    \FOR{each $(x,y)\in\mathcal{D}$}
      \STATE Compute prediction $\hat{y}\gets f_\theta(x)$
      \STATE Record the prediction, reference, and execution cost in $\mathcal{R}$
    \ENDFOR
    \STATE Aggregate metrics over $\mathcal{R}$
    \RETURN Metrics and reproducible evaluation settings
  \end{algorithmic}
\end{algorithm}

\subsection{Reproducibility and limitations}
Report the data split, random seeds, compute settings, and relevant software
versions. State which conclusions are supported by the experiments and
which remain uncertain. Include failure cases and limitations that help
readers understand the method's scope.

\section{Paper publication and open-source workflow}
% Inherited verbatim from the existing team template.
\subsection{Paper publication approval}

Before submitting or publicly releasing a paper, complete the internal publication approval process through the \href{https://yida-group.alibaba-inc.com/APP_ZHVQTG6LH0I0BOSUJDQ8/workbench/FORM-82EC232FEEE9411BB171ECBA40FCD614T9QF?spm=ata.25287382.0.0.27147536laDwbn}{Yida paper publication approval form}. Complete the required review and approval steps before any external publication.

\subsection{Code open-source approval}

If the project will release source code, models, datasets, or related technical assets, initiate and complete the compliance review through the \href{https://lc.alibaba-inc.com/\#/lc/home?survey=1}{Alibaba open-source approval platform} before public release.



\section{Final checklist}
Confirm the title, author order, affiliation, and contact details. Replace
all sample text, figure placeholders, and table dashes. Check references,
figure readability, and PDF metadata. Complete the team's publication and,
where applicable, open-source review before external release.

\begin{ack}
Replace this paragraph with the appropriate acknowledgments and funding
disclosure, or remove the environment when it is not applicable.
\end{ack}

\bibliographystyle{plainnat}
\bibliography{references}

\clearpage
\appendix
\section{Supplementary material}
\label{app:supp}
Use this appendix for additional implementation details, derivations,
evaluation protocols, or qualitative examples. Appendix sections use
lettered headings while equations, tables, and figures retain their normal
counters. Remove this section if no supplementary material is needed.

\subsection{Build and edit guide}
\begin{table}[h]
  \centering
  \caption{Files commonly edited when starting a new report.}
  \begin{tabularx}{\linewidth}{@{}lX@{}}
    \toprule
    File & Purpose\\
    \midrule
    \texttt{yuvion\_team\_neurips.tex} & Title, authors, abstract, and manuscript body.\\
    \texttt{references.bib} & Bibliographic records for cited work.\\
    \texttt{yuvion-neurips.sty} & Team header and technical-report notice.\\
    \texttt{neurips\_2026.sty} & Official NeurIPS style, preserved unchanged.\\
    \bottomrule
  \end{tabularx}
\end{table}

For Overleaf, upload the source ZIP as a new project, choose the main file,
and select pdfLaTeX. For a local TeX Live installation, compile with:
\begin{verbatim}
pdflatex yuvion_team_neurips
bibtex yuvion_team_neurips
pdflatex yuvion_team_neurips
pdflatex yuvion_team_neurips
\end{verbatim}
The second and third LaTeX passes resolve citations and cross-references.
Alternatively, run \verb|latexmk -pdf yuvion_team_neurips.tex|.

\subsection{Conference submission variant}
The default document intentionally displays team identity. To prepare a
conference submission, start from the official author kit, use the required
review mode, remove team branding and identifying material, and add the
official checklist when required. The team header is a report customization,
not part of the official conference submission format.

\end{document}
